\textbf{Phần 1. Trắc nghiệm lựa chọn}
\setcounterpageref{ex}{0}

\begin{ex} %[5D10V57]
Một ô tô xuất phát từ Hà Nội lúc $6\text{ giờ } 45\text{ phút}$, đến Hải Phòng lúc $9\text{ giờ } 10\text{ phút}$. Dọc đường lái xe nghỉ uống nước hết $25\text{ phút}$. Thời gian thực tế ô tô di chuyển trên đường là
\choice
{$1\text{ giờ } 35\text{ phút}$}
{\True $2\text{ giờ}$}
{$2\text{ giờ } 10\text{ phút}$}
{$2\text{ giờ } 25\text{ phút}$}
\loigiai{
Tổng thời gian từ lúc xuất phát đến lúc tới nơi là:
$$9\text{ giờ } 10\text{ phút} - 6\text{ giờ } 45\text{ phút} = 8\text{ giờ } 70\text{ phút} - 6\text{ giờ } 45\text{ phút} = 2\text{ giờ } 25\text{ phút}.$$
Thời gian thực tế xe chạy là:
$$2\text{ giờ } 25\text{ phút} - 25\text{ phút} = 2\text{ giờ}.$$
\textbf{Đáp án của câu hỏi là B}
}
\end{ex}

\begin{ex} %[5D10V57]
Kết quả của phép tính $4\text{ giờ } 35\text{ phút} + 3\text{ giờ } 48\text{ phút} - 2\text{ giờ } 55\text{ phút}$ là
\choice
{$5\text{ giờ } 18\text{ phút}$}
{\True $5\text{ giờ } 28\text{ phút}$}
{$5\text{ giờ } 38\text{ phút}$}
{$4\text{ giờ } 58\text{ phút}$}
\loigiai{
Ta có: $4\text{ giờ } 35\text{ phút} + 3\text{ giờ } 48\text{ phút} = 7\text{ giờ } 83\text{ phút} = 8\text{ giờ } 23\text{ phút}$.\\
Lấy $8\text{ giờ } 23\text{ phút} - 2\text{ giờ } 55\text{ phút} = 7\text{ giờ } 83\text{ phút} - 2\text{ giờ } 55\text{ phút} = 5\text{ giờ } 28\text{ phút}$.\\
\textbf{Đáp án của câu hỏi là B}
}
\end{ex}

\begin{ex} %[5D10V57]
Một người thợ làm việc từ $7\text{ giờ } 30\text{ phút}$ sáng đến $11\text{ giờ } 45\text{ phút}$ trưa. Buổi chiều người đó làm tiếp từ $13\text{ giờ } 15\text{ phút}$ đến $17\text{ giờ}$. Tổng thời gian người đó đã làm việc trong cả ngày là
\choice
{$7\text{ giờ } 30\text{ phút}$}
{\True $8\text{ giờ}$}
{$8\text{ giờ } 15\text{ phút}$}
{$7\text{ giờ } 45\text{ phút}$}
\loigiai{
Thời gian làm buổi sáng:
$$11\text{ giờ } 45\text{ phút} - 7\text{ giờ } 30\text{ phút} = 4\text{ giờ } 15\text{ phút}.$$
Thời gian làm buổi chiều:
$$17\text{ giờ} - 13\text{ giờ } 15\text{ phút} = 16\text{ giờ } 60\text{ phút} - 13\text{ giờ } 15\text{ phút} = 3\text{ giờ } 45\text{ phút}.$$
Tổng thời gian làm việc trong cả ngày là:
$$4\text{ giờ } 15\text{ phút} + 3\text{ giờ } 45\text{ phút} = 7\text{ giờ } 60\text{ phút} = 8\text{ giờ}.$$
\textbf{Đáp án của câu hỏi là B}
}
\end{ex}

\begin{ex} %[5D10V57]
Trong một hội thi bơi lội, vận động viên A bơi hết $2$ phút $15$ giây, vận động viên B bơi nhanh hơn vận động viên A $18$ giây. Vận động viên B bơi hết thời gian là
\choice
{$2\text{ phút } 33\text{ giây}$}
{\True $1\text{ phút } 57\text{ giây}$}
{$1\text{ phút } 47\text{ giây}$}
{$2\text{ phút } 03\text{ giây}$}
\loigiai{
Vì bơi nhanh hơn nên thời gian bơi của vận động viên B ít hơn:\\
$$2\text{ phút } 15\text{ giây} - 18\text{ giây} = 1\text{ phút } 75\text{ giây} - 18\text{ giây} = 1\text{ phút } 57\text{ giây}.$$
\textbf{Đáp án của câu hỏi là B}
}
\end{ex}

\begin{ex} %[5D10V57]
Một đoàn tàu hỏa rời ga A lúc $20\text{ giờ } 30\text{ phút}$ tối thứ Sáu và đến ga B sau $18\text{ giờ } 45\text{ phút}$. Hỏi đoàn tàu đến ga B vào thời điểm nào?
\choice
{$14\text{ giờ } 15\text{ phút}$ ngày thứ Bảy}
{\True $15\text{ giờ } 15\text{ phút}$ ngày thứ Bảy}
{$15\text{ giờ } 45\text{ phút}$ ngày thứ Bảy}
{$16\text{ giờ } 15\text{ phút}$ ngày thứ Bảy}
\loigiai{
Thời gian từ $20\text{ giờ } 30\text{ phút}$ đến nửa đêm ($24\text{ giờ}$) của thứ Sáu là:
$$24\text{ giờ} - 20\text{ giờ } 30\text{ phút} = 3\text{ giờ } 30\text{ phút}.$$
Thời gian tàu chạy tiếp trong ngày thứ Bảy là:
$$18\text{ giờ } 45\text{ phút} - 3\text{ giờ } 30\text{ phút} = 15\text{ giờ } 15\text{ phút}.$$
Vậy đoàn tàu đến ga B lúc $15\text{ giờ } 15\text{ phút}$ ngày thứ Bảy.\\
\textbf{Đáp án của câu hỏi là B}
}
\end{ex}

\begin{ex} %[5D10V57]
Thực hiện phép tính: $5\text{ ngày } 14\text{ giờ} - 2\text{ ngày } 19\text{ giờ}$. Kết quả là
\choice
{$3\text{ ngày } 5\text{ giờ}$}
{\True $2\text{ ngày } 19\text{ giờ}$}
{$2\text{ ngày } 17\text{ giờ}$}
{$3\text{ ngày } 19\text{ giờ}$}
\loigiai{
Đổi: $5\text{ ngày } 14\text{ giờ} = 4\text{ ngày } 38\text{ giờ}$.\\
Thực hiện phép trừ:
$$4\text{ ngày } 38\text{ giờ} - 2\text{ ngày } 19\text{ giờ} = 2\text{ ngày } 19\text{ giờ}.$$
\textbf{Đáp án của câu hỏi là B}
}
\end{ex}

\begin{ex} %[5D10V57]
Bác Bình may xong một chiếc áo hết $1\text{ giờ } 25\text{ phút}$, may xong một chiếc quần hết ít hơn may chiếc áo $15\text{ phút}$. Hỏi bác Bình may xong cả bộ quần áo đó hết bao nhiêu thời gian?
\choice
{$2\text{ giờ } 25\text{ phút}$}
{\True $2\text{ giờ } 35\text{ phút}$}
{$2\text{ giờ } 40\text{ phút}$}
{$2\text{ giờ } 15\text{ phút}$}
\loigiai{
Thời gian may một chiếc quần là:
$$1\text{ giờ } 25\text{ phút} - 15\text{ phút} = 1\text{ giờ } 10\text{ phút}.$$
Thời gian may cả bộ quần áo là:
$$1\text{ giờ } 25\text{ phút} + 1\text{ giờ } 10\text{ phút} = 2\text{ giờ } 35\text{ phút}.$$
\textbf{Đáp án của câu hỏi là B}
}
\end{ex}

\begin{ex} %[5D10V57]
Một ca nô đi xuôi dòng từ bến A đến bến B hết $1\text{ giờ } 45\text{ phút}$. Khi đi ngược dòng từ B về A, do nước cản nên ca nô đi lâu hơn xuôi dòng $30\text{ phút}$. Tổng thời gian ca nô cả đi lẫn về (không tính thời gian nghỉ tại B) là
\choice
{$3\text{ giờ } 45\text{ phút}$}
{\True $4\text{ giờ}$}
{$4\text{ giờ } 15\text{ phút}$}
{$3\text{ giờ } 15\text{ phút}$}
\loigiai{
Thời gian ca nô đi ngược dòng từ B về A là:
$$1\text{ giờ } 45\text{ phút} + 30\text{ phút} = 2\text{ giờ } 15\text{ phút}.$$
Tổng thời gian ca nô đi cả hai chiều là:
$$1\text{ giờ } 45\text{ phút} + 2\text{ giờ } 15\text{ phút} = 3\text{ giờ } 60\text{ phút} = 4\text{ giờ}.$$
\textbf{Đáp án của câu hỏi là B}
}
\end{ex}

\begin{ex} %[5D10V57]
Một học sinh làm xong bài tập Toán hết $35$ phút, làm bài tập Tiếng Việt hết $40$ phút. Thời gian bạn đó làm cả hai bài tập là
\choice
{$1\text{ giờ } 10\text{ phút}$}
{\True $1\text{ giờ } 15\text{ phút}$}
{$1\text{ giờ } 20\text{ phút}$}
{$1\text{ giờ } 25\text{ phút}$}
\loigiai{
Tổng thời gian làm hai bài tập là:
$$35\text{ phút} + 40\text{ phút} = 75\text{ phút} = 1\text{ giờ } 15\text{ phút}.$$
\textbf{Đáp án của câu hỏi là B}
}
\end{ex}

\begin{ex} %[5D10V57]
Một trận đấu bóng đá bắt đầu lúc $15\text{ giờ } 30\text{ phút}$. Thời gian thi đấu gồm $2$ hiệp, mỗi hiệp $45$ phút và nghỉ giữa hai hiệp $15$ phút. Trận đấu bóng đá đó kết thúc lúc
\choice
{$17\text{ giờ}$}
{\True $17\text{ giờ } 15\text{ phút}$}
{$17\text{ giờ } 30\text{ phút}$}
{$17\text{ giờ } 45\text{ phút}$}
\loigiai{
Tổng thời gian thi đấu và nghỉ giữa hai hiệp là:
$$45\text{ phút} \times 2 + 15\text{ phút} = 90\text{ phút} + 15\text{ phút} = 105\text{ phút} = 1\text{ giờ } 45\text{ phút}.$$
Trận đấu kết thúc vào lúc:
$$15\text{ giờ } 30\text{ phút} + 1\text{ giờ } 45\text{ phút} = 16\text{ giờ } 75\text{ phút} = 17\text{ giờ } 15\text{ phút}.$$
\textbf{Đáp án của câu hỏi là B}
}
\end{ex}

\begin{ex} %[5D10V57]
Một gia đình đi xe máy về quê phải mất $2\text{ giờ } 20\text{ phút}$. Họ muốn đến quê lúc $10\text{ giờ}$ sáng và dự định nghỉ dọc đường $15\text{ phút}$. Hỏi gia đình đó phải xuất phát muộn nhất lúc mấy giờ?
\choice
{$7\text{ giờ } 15\text{ phút}$}
{\True $7\text{ giờ } 25\text{ phút}$}
{$7\text{ giờ } 35\text{ phút}$}
{$7\text{ giờ } 45\text{ phút}$}
\loigiai{
Tổng thời gian đi và nghỉ là:
$$2\text{ giờ } 20\text{ phút} + 15\text{ phút} = 2\text{ giờ } 35\text{ phút}.$$
Thời điểm xuất phát muộn nhất là:
$$10\text{ giờ} - 2\text{ giờ } 35\text{ phút} = 9\text{ giờ } 60\text{ phút} - 2\text{ giờ } 35\text{ phút} = 7\text{ giờ } 25\text{ phút}.$$
\textbf{Đáp án của câu hỏi là B}
}
\end{ex}

\begin{ex} %[5D10V57]
Kết quả của phép tính: $12\text{ phút } 10\text{ giây} - 7\text{ phút } 45\text{ giây} + 3\text{ phút } 50\text{ giây}$ là
\choice
{$7\text{ phút } 15\text{ giây}$}
{\True $8\text{ phút } 15\text{ giây}$}
{$8\text{ phút } 25\text{ giây}$}
{$7\text{ phút } 55\text{ giây}$}
\loigiai{
Thực hiện phép trừ trước:
$$12\text{ phút } 10\text{ giây} - 7\text{ phút } 45\text{ giây} = 11\text{ phút } 70\text{ giây} - 7\text{ phút } 45\text{ giây} = 4\text{ phút } 25\text{ giây}.$$
Tiếp tục cộng:
$$4\text{ phút } 25\text{ giây} + 3\text{ phút } 50\text{ giây} = 7\text{ phút } 75\text{ giây} = 8\text{ phút } 15\text{ giây}.$$
\textbf{Đáp án của câu hỏi là B}
}
\end{ex}

\textbf{Phần 2. Trắc nghiệm đúng sai}
\setcounterpageref{ex}{0}

\begin{ex} %[5D10V57]
Cho các phép tính cộng, trừ số đo thời gian sau:
\choiceTF
{\True $2\text{ giờ } 45\text{ phút} + 3\text{ giờ } 25\text{ phút} = 6\text{ giờ } 10\text{ phút}$}
{\True $5\text{ ngày } 6\text{ giờ} - 1\text{ ngày } 18\text{ giờ} = 3\text{ ngày } 12\text{ giờ}$}
{$4\text{ phút } 12\text{ giây} - 2\text{ phút } 30\text{ giây} = 1\text{ phút } 18\text{ giây}$}
{\True $7\text{ giờ } 20\text{ phút} - 4\text{ giờ } 50\text{ phút} = 2\text{ giờ } 30\text{ phút}$}
\loigiai{
\begin{itemchoice}
\item Đúng. $2\text{ giờ } 45\text{ phút} + 3\text{ giờ } 25\text{ phút} = 5\text{ giờ } 70\text{ phút} = 6\text{ giờ } 10\text{ phút}$.
\item Đúng. $5\text{ ngày } 6\text{ giờ} = 4\text{ ngày } 30\text{ giờ} \Rightarrow 4\text{ ngày } 30\text{ giờ} - 1\text{ ngày } 18\text{ giờ} = 3\text{ ngày } 12\text{ giờ}$.
\item Sai. $4\text{ phút } 12\text{ giây} = 3\text{ phút } 72\text{ giây} \Rightarrow 3\text{ phút } 72\text{ giây} - 2\text{ phút } 30\text{ giây} = 1\text{ phút } 42\text{ giây}$ (chứ không phải $1\text{ phút } 18\text{ giây}$).
\item Đúng. $7\text{ giờ } 20\text{ phút} = 6\text{ giờ } 80\text{ phút} \Rightarrow 6\text{ giờ } 80\text{ phút} - 4\text{ giờ } 50\text{ phút} = 2\text{ giờ } 30\text{ phút}$.
\end{itemchoice}
}
\end{ex}

\begin{ex} %[5D10V57]
Ba bạn An, Bình, Cường cùng tham gia một cuộc thi chạy $800\text{ m}$. An chạy hết $2\text{ phút } 45\text{ giây}$, Bình chạy nhiều hơn An $15\text{ giây}$, Cường chạy ít hơn Bình $20\text{ giây}$. Xét tính đúng/sai của các nhận định sau:
\choiceTF
{\True Bình chạy hết $3\text{ phút}$}
{\True Cường chạy hết $2\text{ phút } 40\text{ giây}$}
{Bình là người về đích đầu tiên}
{\True Cường là người chạy nhanh nhất và về đích đầu tiên}
\loigiai{
\begin{itemchoice}
\item Đúng. Thời gian của Bình: $2\text{ phút } 45\text{ giây} + 15\text{ giây} = 3\text{ phút}$.
\item Đúng. Thời gian của Cường: $3\text{ phút} - 20\text{ giây} = 2\text{ phút } 40\text{ giây}$.
\item Sai. Người có thời gian chạy ít nhất là người về đích đầu tiên. Bình chạy mất $3$ phút (nhiều thời gian nhất) nên về đích cuối cùng.
\item Đúng. Cường mất $2\text{ phút } 40\text{ giây}$ (ít thời gian nhất) nên Cường về đích đầu tiên.
\end{itemchoice}
}
\end{ex}

\textbf{Phần 3. Câu trả lời ngắn}
\setcounterpageref{ex}{0}

\begin{ex} %[5D10V57]
Một người đi từ tỉnh A lúc $6\text{ giờ } 15\text{ phút}$ và đến tỉnh B lúc $10\text{ giờ } 45\text{ phút}$. Dọc đường người đó dừng lại ăn sáng hết $30\text{ phút}$. Hỏi thời gian thực tế người đó đi trên đường là bao nhiêu phút?
\shortans{240}
\loigiai{
Tổng thời gian từ lúc xuất phát đến lúc tới nơi là:
$$10\text{ giờ } 45\text{ phút} - 6\text{ giờ } 15\text{ phút} = 4\text{ giờ } 30\text{ phút}.$$
Thời gian thực tế đi trên đường là:
$$4\text{ giờ } 30\text{ phút} - 30\text{ phút} = 4\text{ giờ}.$$
Đổi $4\text{ giờ} = 240\text{ phút}$.
}
\end{ex}

\begin{ex} %[5D10V57]
Bác thợ hàn hoàn thành khung cửa thứ nhất hết $2$ giờ $15$ phút, khung cửa thứ hai làm hết nhiều hơn khung thứ nhất $25$ phút. Hỏi bác hoàn thành cả hai khung cửa hết bao nhiêu phút?
\shortans{295}
\loigiai{
Thời gian làm khung cửa thứ hai là:
$$2\text{ giờ } 15\text{ phút} + 25\text{ phút} = 2\text{ giờ } 40\text{ phút}.$$
Tổng thời gian làm cả hai khung cửa là:
$$2\text{ giờ } 15\text{ phút} + 2\text{ giờ } 40\text{ phút} = 4\text{ giờ } 55\text{ phút}.$$
Đổi: $4\text{ giờ } 55\text{ phút} = 4 \times 60 + 55 = 295\text{ (phút)}$.
}
\end{ex}

\begin{ex} %[5D10V57]
Một vận động viên hoàn thành chặng 1 giải chạy trong $48$ phút $35$ giây, hoàn thành chặng 2 trong $52$ phút $40$ giây. Tổng thời gian hoàn thành cả hai chặng của vận động viên đó là bao nhiêu giây?
\shortans{6075}
\loigiai{
Tổng thời gian là:
$$48\text{ phút } 35\text{ giây} + 52\text{ phút } 40\text{ giây} = 100\text{ phút } 75\text{ giây} = 101\text{ phút } 15\text{ giây}.$$
Đổi ra giây:
$$101 \times 60 + 15 = 6\,060 + 15 = 6\,075\text{ (giây)}.$$
}
\end{ex}

\begin{ex} %[5D10V57]
Hôm nay bạn Lan học ở trường từ $7\text{ giờ } 15\text{ phút}$ đến $11\text{ giờ } 15\text{ phút}$. Buổi chiều Lan học thêm từ $14\text{ giờ}$ đến $16\text{ giờ } 30\text{ phút}$. Hỏi hôm nay Lan đã học ở trường tất cả bao nhiêu phút?
\shortans{390}
\loigiai{
Thời gian học buổi sáng:
$$11\text{ giờ } 15\text{ phút} - 7\text{ giờ } 15\text{ phút} = 4\text{ giờ}.$$
Thời gian học buổi chiều:
$$16\text{ giờ } 30\text{ phút} - 14\text{ giờ} = 2\text{ giờ } 30\text{ phút}.$$
Tổng thời gian học là:
$$4\text{ giờ} + 2\text{ giờ } 30\text{ phút} = 6\text{ giờ } 30\text{ phút}.$$
Đổi $6\text{ giờ } 30\text{ phút} = 6 \times 60 + 30 = 390\text{ (phút)}$.
}
\end{ex}

\begin{ex} %[5D10V57]
Một người bắt đầu tập thể dục lúc $5\text{ giờ } 45\text{ phút}$ sáng và kết thúc lúc $6\text{ giờ } 35\text{ phút}$ sáng. Hỏi người đó đã tập thể dục trong bao nhiêu phút?
\shortans{50}
\loigiai{
Thời gian tập thể dục là:
$$6\text{ giờ } 35\text{ phút} - 5\text{ giờ } 45\text{ phút} = 5\text{ giờ } 95\text{ phút} - 5\text{ giờ } 45\text{ phút} = 50\text{ (phút)}.$$
}
\end{ex}

\begin{ex} %[5D10V57]
Một chuyến xe khách xuất phát lúc $13\text{ giờ } 45\text{ phút}$ và đi hết $3\text{ giờ } 50\text{ phút}$. Xe khách tới nơi lúc mấy giờ mấy phút? (Nhập số phút của giờ đến nơi, ví dụ nếu đến nơi lúc 17 giờ 35 phút thì kết quả là 35).
\shortans{35}
\loigiai{
Thời điểm xe khách tới nơi là:
$$13\text{ giờ } 45\text{ phút} + 3\text{ giờ } 50\text{ phút} = 16\text{ giờ } 95\text{ phút} = 17\text{ giờ } 35\text{ phút}.$$
Số phút là $35$.
}
\end{ex}
